\documentclass[11pt]{article}
\usepackage[a4paper,margin=22mm]{geometry}
\usepackage{fontspec}
\setmainfont{DejaVu Serif}
\setsansfont{DejaVu Sans}
\usepackage{microtype}
\usepackage{xcolor}
\usepackage{hyperref}
\usepackage{enumitem}
\usepackage{parskip}
\definecolor{Accent}{HTML}{971D32}
\definecolor{Muted}{HTML}{666666}
\hypersetup{colorlinks=true,urlcolor=Accent,linkcolor=Accent}
\setlist{leftmargin=1.6em,itemsep=0.3em,topsep=0.35em}
\setcounter{tocdepth}{2}
\renewcommand{\contentsname}{Contents}
\newcommand{\sectionrule}{\par\vspace{-0.5em}\noindent\textcolor{Accent}{\rule{\linewidth}{0.6pt}}\par}
\begin{document}

\begin{center}
{\sffamily\bfseries\color{Accent} TOEFL WORD OF THE DAY\par}
\vspace{8mm}
{\fontsize{42}{48}\selectfont\bfseries innovative\par}
\vspace{2mm}
{\Large\sffamily\color{Accent} /ˈɪnəveɪtɪv/\par}
\vspace{2mm}
{\small\sffamily Local audio phrase: innovative\par}
{\small\sffamily Audio file: audios/innovative.mp3\par}
\vspace{2mm}
{\large\sffamily\color{Muted} adjective\par}
\vspace{5mm}
{\sffamily introducing useful new ideas or methods\par}
\vspace{6mm}
{\large Innovative describes something that introduces a genuinely new or imaginative way of doing things.\par}
\vspace{6mm}
\begin{minipage}{0.84\textwidth}
\itshape “The researchers developed an innovative method for removing pollutants from drinking water.”
\end{minipage}\par
\vspace{10mm}
\textbf{Date:} September 15th 2026\quad\textbullet\quad \textbf{Level:} B1--mid-B2 TOEFL
\vfill
Created and compiled by \textbf{Amir Kooshky}\par
\vspace{2mm}
\href{https://t.me/KooshkyTOEFL}{Telegram Channel}\quad\textbullet\quad
\href{https://www.instagram.com/kooshkytoefl}{Instagram}\quad\textbullet\quad
\href{https://t.me/KooshkyTOEFL\_pv}{Personal Telegram}
\end{center}
\newpage
\tableofcontents
\newpage

\section{Meanings and usage}
\sectionrule
\subsection{Meaning 1: Introducing useful new ideas or methods}
\textbf{Part of speech:} adjective

\textbf{IPA:} /ˈɪnəveɪtɪv/

\textbf{Pronunciation phrase:} innovative

\textbf{Local audio file:} audios/innovative.mp3

\textbf{Register:} neutral to formal; common in academic, business, technology, and policy contexts

\textbf{Definition:} introducing or using new ideas, methods, or designs in an original and useful way

\textbf{In simpler words:} new, original, and useful

\textbf{Typical pattern:} an innovative + approach, method, solution, design, product, or idea

\subsubsection{Natural examples}
\begin{enumerate}
  \item The researchers developed an innovative method for removing pollutants from drinking water.
  \item The city adopted an innovative approach to reducing traffic in crowded neighborhoods.
  \item Innovative teaching practices can encourage students to participate more actively.
  \item The company is known for innovative products that use less energy.
  \item Engineers proposed an innovative solution to the region's water shortage.
  \item The museum's innovative use of digital displays made the exhibition more accessible.
  \item Small firms may struggle to finance innovative projects despite their potential value.
  \item Her innovative research combined methods from psychology and computer science.
\end{enumerate}

\subsubsection{Nuance and implication}
\textbf{Central nuance:} Innovative suggests more than being recent: it introduces an original improvement or a different way of working.

\textbf{Usual implication:} The new idea or method has practical value or the potential to improve existing practice.

\textbf{Compared with new:} Anything recently created can be new. Innovative is used when it also shows originality or introduces a meaningful change.

\subsubsection{Grammar notes}
\textbf{How to use it:} Use it before a noun, as in an innovative method, or after a linking verb, as in The design is innovative.

\textbf{What can follow it:} It commonly describes an approach, method, technique, solution, idea, design, product, program, company, or researcher.

\textbf{Position in a sentence:} Words such as highly, particularly, genuinely, and technologically can appear before innovative.

\subsubsection{Useful patterns}
\textbf{Pattern:} an innovative approach to + noun or -ing form

\textbf{Use:} Use this for an original way of dealing with a subject or problem.

\textbf{Example:} The course takes an innovative approach to teaching academic vocabulary.

\textbf{Pattern:} an innovative solution to + problem

\textbf{Use:} Use this for an original and practical answer to a difficulty.

\textbf{Example:} The team proposed an innovative solution to food waste.

\textbf{Pattern:} innovative use of + noun

\textbf{Use:} Use this when something familiar is applied in a new way.

\textbf{Example:} The project demonstrates an innovative use of recycled materials.

\subsubsection{Common wrong usage}
\paragraph{Correction 1}
\textbf{Wrong:} The ancient tool was innovative because it was very old.

\textbf{Correct:} The ancient tool was historically important because it was very old.

\textbf{Why:} Innovative refers to originality or the introduction of a new method, not simply to age or importance.

\paragraph{Correction 2}
\textbf{Wrong:} They developed an innovation solution.

\textbf{Correct:} They developed an innovative solution.

\textbf{Why:} Use the adjective innovative before a noun. Innovation is a noun.

\paragraph{Correction 3}
\textbf{Wrong:} The laboratory innovated a new testing method.

\textbf{Correct:} The laboratory developed an innovative testing method.

\textbf{Why:} Innovate is usually used without a direct object; develop an innovative method is more natural here.

\section{Word family}
\sectionrule
\subsection{innovation}
\textbf{Part of speech:} countable or uncountable noun

\textbf{IPA:} /ˌɪnəˈveɪʃən/

\textbf{Definition:} a new idea, method, or product, or the process of introducing such changes

\textbf{Relationship to the headword:} This noun names an innovative idea or the process of creating and introducing new ideas.

\textbf{Grammar pattern:} innovation in + field, or an innovation in + product or method

\textbf{Usage note:} It is uncountable for the general process and countable for a particular new development.

\subsubsection{Examples}
\begin{enumerate}
  \item Technological innovation has transformed access to information.
  \item The device was an important innovation in medical imaging.
\end{enumerate}

\subsection{innovate}
\textbf{Part of speech:} intransitive verb

\textbf{IPA:} /ˈɪnəveɪt/

\textbf{Definition:} to introduce new ideas, methods, or products

\textbf{Relationship to the headword:} This verb describes the action of creating or introducing innovations.

\textbf{Grammar pattern:} innovate in + field, or innovate by + -ing form

\textbf{Usage note:} It is usually used without a direct object.

\subsubsection{Examples}
\begin{enumerate}
  \item Organizations must innovate to respond to changing needs.
  \item The company continued to innovate despite limited resources.
\end{enumerate}

\subsection{innovator}
\textbf{Part of speech:} countable noun

\textbf{IPA:} /ˈɪnəveɪtər/

\textbf{Definition:} a person or organization that introduces new ideas or methods

\textbf{Relationship to the headword:} This noun names the person or organization responsible for innovative work.

\subsubsection{Examples}
\begin{enumerate}
  \item The scientist was an early innovator in renewable-energy research.
  \item Successful innovators often test several versions of an idea.
\end{enumerate}

\section{Common collocations}
\sectionrule
\subsection{innovative approach}
\textbf{Applies to:} Introducing useful new ideas or methods

\textbf{Meaning:} an original way of dealing with a task or subject

\textbf{Example:} The study uses an innovative approach to measuring social trust.

\subsection{innovative solution}
\textbf{Applies to:} Introducing useful new ideas or methods

\textbf{Meaning:} an original and useful answer to a problem

\textbf{Example:} Local engineers found an innovative solution to seasonal flooding.

\subsection{innovative technology}
\textbf{Applies to:} Introducing useful new ideas or methods

\textbf{Meaning:} technology that introduces an important new design or method

\textbf{Example:} Innovative technology may reduce the cost of storing solar energy.

\subsection{highly innovative}
\textbf{Applies to:} Introducing useful new ideas or methods

\textbf{Meaning:} showing a great degree of originality

\textbf{Example:} At the time, the treatment was considered highly innovative.

\section{Synonyms and antonyms}
\sectionrule
\subsection{Close synonyms}
\subsubsection{original}
\textbf{Part of speech:} adjective

\textbf{Applies to:} Introducing useful new ideas or methods

\textbf{Shared meaning:} Both can describe ideas that are new and creative.

\textbf{Difference:} Original emphasizes that an idea is new and not copied. Innovative also suggests applying originality to improve how something is done.

\textbf{Example:} The paper offers an original interpretation of the historical evidence.

\subsubsection{inventive}
\textbf{Part of speech:} adjective

\textbf{Applies to:} Introducing useful new ideas or methods

\textbf{Shared meaning:} Both praise creative problem-solving.

\textbf{Difference:} Inventive emphasizes imagination and skill in creating things. Innovative often emphasizes introducing a new approach in practice.

\textbf{Example:} The engineers devised an inventive way to reuse waste heat.

\subsection{Direct or useful antonyms}
\subsubsection{conventional}
\textbf{Part of speech:} adjective

\textbf{Applies to:} Introducing useful new ideas or methods

\textbf{Opposition:} Conventional methods follow established practice, whereas innovative methods introduce original changes.

\textbf{Example:} The control group received conventional treatment.

\section{Words learners may confuse}
\sectionrule
\subsection{innovation}
\textbf{Part of speech:} adjective versus noun

\textbf{Why learners confuse them:} Innovative and innovation belong to the same word family but occupy different positions in a sentence.

\textbf{Key difference:} Innovative is an adjective; innovation is a noun.

\textbf{innovative:} The team tested an innovative design.

\textbf{innovation:} The design represents an important innovation.

\section{Etymology}
\sectionrule
\textbf{Origin:} Innovative is formed from innovate, which comes from Latin innovare, meaning to renew or make new.

\textbf{Memory link:} The nov element also appears in novel and connects innovative with the idea of newness.

\section{Practice exercises}
\sectionrule
\subsection{Word family choice}
Choose the best member of the word family for each blank.

\begin{enumerate}
  \item The device was a major \_\_\_\_\_ in low-cost water purification.
  \item Organizations need to \_\_\_\_\_ when existing methods stop working.
  \item The architect proposed an \_\_\_\_\_ use of the narrow space.
\end{enumerate}

\subsection{Correct or incorrect?}
Decide whether each sentence is correct. Correct the inaccurate ones.

\begin{enumerate}
  \item The researchers developed an innovative method for analyzing the samples.
  \item The team created an innovation approach.
  \item The product is new, but its design is not especially innovative.
\end{enumerate}

\subsection{Collocation practice}
Complete each sentence with a natural collocation.

\begin{enumerate}
  \item The city adopted an innovative \_\_\_\_\_ to managing traffic.
  \item The engineers proposed an innovative \_\_\_\_\_ to the storage problem.
\end{enumerate}

\subsection{Complete answer key}
\subsubsection{Word family choice}
\begin{enumerate}
  \item \textbf{innovation.} The device was a major innovation in low-cost water purification. \emph{Use the countable noun innovation after a major.}
  \item \textbf{innovate.} Organizations need to innovate when existing methods stop working. \emph{Use the verb innovate after need to.}
  \item \textbf{innovative.} The architect proposed an innovative use of the narrow space. \emph{Use the adjective innovative before use.}
\end{enumerate}

\subsubsection{Correct or incorrect?}
\begin{enumerate}
  \item \textbf{Correct} \emph{Innovative method is a natural collocation.}
  \item \textbf{Incorrect}. The team created an innovative approach. \emph{Use the adjective innovative before approach.}
  \item \textbf{Correct} \emph{Something can be recently made without introducing an original improvement.}
\end{enumerate}

\subsubsection{Collocation practice}
\begin{enumerate}
  \item \textbf{approach.} The city adopted an innovative approach to managing traffic. \emph{Innovative approach to is an established pattern.}
  \item \textbf{solution.} The engineers proposed an innovative solution to the storage problem. \emph{Innovative solution to is a natural collocation.}
\end{enumerate}

\vfill
\begin{center}
\small\color{Muted} Created and compiled by \textbf{Amir Kooshky}\\
\href{https://t.me/KooshkyTOEFL}{Telegram Channel}\quad\textbullet\quad
\href{https://www.instagram.com/kooshkytoefl}{Instagram}\quad\textbullet\quad
\href{https://t.me/KooshkyTOEFL\_pv}{Personal Telegram}
\end{center}
\end{document}
